\recipe{Recipe Template}
\label{recipe:recipe_template}
\makes{x portions}
\worktime{\SI{10} \minute}
\totaltime{\SI{10} \minute}
\source{\href{https://www.bbcgoodfood.com/}{short link to source}}

\begin{ingredients}

  \item Udon noodles, 2 packs
  \item Tofu or chicken, \SI{200}{\gram}
  \item Water/Stock, \SI{2.5}{\litre}
  \item Coconut milk, \SI{250}{\milli \litre}
  \item Garlic, \SI{1}{\tablespoon} crushed
  \item Turmeric, \SI{1.5}{\teaspoon}

  \ingroup{To Serve}
  \item 2 spring onions, chopped
  \item Naan bread or rice
\end{ingredients}

\begin{method}
  	A recipe must fit onto one page
  
  	A recipe must have an image, preferably one taken by the creator with friends.

  	Include amounts in the method, such as fry for \SI{30}{\second}.

\end{method}

\begin{figure}[h]
  	\centering
  	
  	\begin{tikzpicture}
  		% Save the image in a box to get its dimensions
  		\sbox0{\includegraphics[width = 0.6\linewidth]{images/MissingImage.png}}%
  
  		% Draw a clipping rectangle with rounded corners
  		\path[clip, rounded corners=0.5cm] (0,0) rectangle (\wd0,\ht0);
  
  		% Place the image inside the clipped area
  		\node[inner sep=0pt, outer sep=0pt, anchor=south west] at (0,0) {\usebox0};
	\end{tikzpicture}

  \caption*{caption for image with date and information, 1st January 2026}
\end{figure}

\begin{notes}
notes on recipe go here.
\end{notes}

